In this section I will show a synthetic example that helps explaining why the problem is relevant and shows the difference perfomance between naive methods and the semiparametric estimation proposed by \citet{mccartan2025identificationsemiparametricestimationconditional} and embedded in R package \texttt{seine} (\citet{seine}).

In the synthetic example we are the ones who generate the data and therefore we know the exact model of the reality.
We show that, if we generate the data and then we try to estimate with different methods we obtain different results and the naive methods do not work well.

\subsection{Italian constitutional referendum: a motivating example}

In particular, we can exploit the constitutional referendum in Italy to show a possible way in which the naive estimators could not work and therefore we can generate a dataset that encompasses the problem.

The problem is the following: we want to estimate how did voters with different levels of income vote.

If we look at data at regional level, we would be tempted to say that richest people voted more in favour of the government proposal. In fact, the only three regions where the votes in favour were more than the votes in contrast (Lombardy, Veneto and Friuli-Venezia Giulia) are among the richest regions in Italy. On the contrary, the regions were No got more than 60\% are among the poorest regions of Italy (Basilicata, Campania, Sicily).

However, once we look at data at municipal level, we find that, within the regions, the relationship is nearly inverse: No votes were more frequent in urban areas and other municipalities where income is relatively higher.

That divergence makes us think that if we were able to observe results at individual level, we could find even different relationships between income and voting behavior.

Therefore, we can now generate a synthetic dataset that presents a structure that resembles some stylized facts that we observe in the constitutional referendum results and other notions that we have about electoral behaviors in Italy.

\subsection{The synthetic experiment}

The synthetic experiment is constructed as follows.

We assume that election takes place in a Country which has three different regions: North, Centre and South.

Each region is divided in 100 precincts, so that in the end we observe 300 precincts: 100 precincts in the North region, 100 in the Centre region, 100 in the South region.

Each precinct has the same voting population $N_g$, which is then normalised to 1.

For each electoral precinct, we observe the percentage of high-income and low-income voters ($\overline{X}_g$) and the electoral results ($Y_g$).

Within each income category, there is high variability in their electoral behaviour across different precincts.

In particular, we impose that high-income people in the North and South precincts have a high propensity of voting in favour of the reform, while in the Centre this percentage is much lower.
As for low-income people, we impose that in North they have high propensity of voting in favour, while in the Centre and in the South this propensity is much lower.

At the same time, we impose that the percentage of high-income voters is highest in the North and lowest in the South and intermediate in the Centre.

\subsubsection{Data generating process}

We generate the data as follows. 

For each region $r \in \{\text{North}, \text{Centre}, \text{South}\}$ and for each of the 100 precincts in that region, we draw the share of high-income voters from a truncated normal distribution,

\begin{equation*}
\overline{X}_{gr}^{H} \sim \text{TruncNormal}(a=0, b=1, \mu_r, \sigma=0.2),
\end{equation*}

with means $\mu_{\text{North}}=0.7$, $\mu_{\text{Centre}}=0.5$, and $\mu_{\text{South}}=0.3$. The low-income share is then set to $\overline{X}_{gr}^{L}=1-\overline{X}_{gr}^{H}$. Hence, on average, Northern precincts have higher shares of high-income voters, Southern precincts have lower shares, and Central precincts lie in between.
The use of Truncated Normal distribution allows to have quite high variability within geography ($\sigma = 0.2$), but without having negative shares or shares above 1.

Subsequently, I assign to each precinct a propensity to vote in favour of the reform within each income group and I encode it in a vector $ \beta_g = \begin{pmatrix}
\beta_g^H \\
\beta_g^L
\end{pmatrix}$.

$\beta_g^H$ and $\beta_g^L$ are extracted again from truncated normal distributions with different means depending on the region.

\begin{align*}
\beta_{gr}^{H} &\sim \text{TruncNormal}(a=0, b=1, \mu=\eta_{gr}^{H}, \sigma=0.05), \\
\beta_{gr}^{L} &\sim \text{TruncNormal}(a=0, b=1, \mu=\eta_{gr}^{L}, \sigma=0.05).
\end{align*}

with means $\eta_{gr}^{H}$ and $\eta_{gr}^{L}$ defined as follows:
\begin{align*}
\eta_{\text{North}}^{H} = 0.6, \quad & \eta_{\text{North}}^{L} = 0.7, \\
\eta_{\text{Centre}}^{H} = 0.3, \quad & \eta_{\text{Centre}}^{L} = 0.3, \\
\eta_{\text{South}}^{H} = 0.6, \quad & \eta_{\text{South}}^{L} = 0.2.
\end{align*}

The key feature of this design is that the income-specific propensities are heterogeneous within each group and also vary across regions. High-income voters in the North and South have a strong tendency to support the reform, while in the Centre the same group is much less favourable. Low-income voters have a relatively high propensity to support the reform in the North, but not in the Centre or South.

Finally, the observed precinct-level outcome is generated as the weighted average of these two group-specific propensities:

\begin{equation*}
Y_{g}=\overline{X}_{g}^{H}\beta_{g}^{H}+\overline{X}_{g}^{L}\beta_{g}^{L}.
\end{equation*}

In the end, we have a dataset that for each electoral precinct $g$ contains the following information: the share of high-income voters $\overline{X}_{g}^{H}$, the share of low-income voters $\overline{X}_{g}^{L}$, the observed outcome $Y_{g}$, the location of the precinct (North, Centre or South), the expected propensity to vote in favour of the reform for high-income voters $\eta_{g}^{H}$ and for low-income voters $\eta_{g}^{L}$ (equal across geographies of the same region), and the actual propensities $\beta_{g}^{H}$ and $\beta_{g}^{L}$ (specific of each geography).

\subsubsection{From DGP to the model}

The generated dataset strictly satisfies the assumptions stated in previous Sections.

The response variable $Y_i$ is a binary variable $Y_i \in \{0,1\}$ and it is observed only at geography level as a share $\overline{Y}_g \in [0,1]$.

Population is partitioned across geographies and within each geography, it is partitioned into two groups (high-income and low-income voters), so that for each geography $g$ we observe $\overline{X}_g = \begin{pmatrix}
    \overline{X}_g^H \\
    \overline{X}_g^L
\end{pmatrix}$ with $\overline{X}_g^L = 1- \overline{X}_g^H$.

Moreover, we can interpret the region of each precinct as a covariate $Z_g = \{0,1\}^3$, which in this case does not arise from individual characteristics.

$\overline{Y}_g$, $\overline{X}_g$ and $Z_g$ are observed, while $\beta_g$ (and their corresponding $\eta_g$) are unobserved and are the goals of the semiparametric estimation.

Bounded N is satisfied as $N_g$ is equal across different geographies.

Positivity is satisfied as we have allowed for variability in the shares of high-income and low-income voters across geographies within each region.

Finally, CAR-IND is almost completely satisfied, thanks to the fact that $\beta_g$ are generated from a distribution with low variance ($\sigma = 0.05$) that depends on the covariate $Z_g$.

In fact, a complete adherence to CAR-IND would require $\beta_g$ to be equal across geographies of the same region, but I have allowed for some variability in order to make the example more realistic and more responding to the expected ability of the covariates to explain the variation in the propensitities.

In our case, $E[Y_i | G_i = g, X_i = x_i, Z_{G_i} = z_g ]  = \beta_g$ is just slightly different from $E[Y_i | X_i = x_i, Z_{G_i} = z_g ] = \eta_{gr}$.

\subsection{Comparison of results}

\subsubsection{OLS estimates}

As previously mentioned, in our dataset we know the true values of $\beta_g$ for each geography $g$, therefore we are able to obtain the true values of the estimands of interest $\beta^H$ and $\beta^L$ by aggregating the true values of $\beta_g$ as stated in the first accounting identity in Section 2: $\beta^H = \frac{\sum_{g=1}^G \beta^H_g*N^H_g}{\sum_g N^H_g}$ and $\beta^L = \frac{\sum_{g=1}^G \beta_g*N_g^L}{\sum_g N_g^L}$.

In our case, we have that $\beta_H = 0.503$ and $\beta_L = 0.343$.

The naive OLS specification, implicitly assuming CCAR, would rely on the following regression: $\overline{Y_g} = \beta_H*\overline{X}_g^H + \beta_L*\overline{X}_g^L+\epsilon_g$.
As $\overline{X}_g^H$ and $\overline{X_g^L}$ are perfectly collinear($\overline{X}_g^H + \overline{X}_g^L = 1$), we use $\overline{X}_g^L$ as reference category and we estimate the following regression: 
\begin{align*}
\overline{Y_g} & = \beta_H*\overline{X}_g^H + \beta_L*\overline{X}_g^L+\epsilon_g 
\\& = \beta_H*\overline{X}_g^H + \beta_L*(1-\overline{X}_g^H)+\epsilon_g
\\& = \beta_L + (\beta_H - \beta_L)*\overline{X}_g^H + \epsilon_g.
\end{align*}

Results are presented in Tab. \ref{tab:naive-regression} and show that $\beta_L$ is estimated to be 0.213, while $\beta_H$ is estimated to be 0.213+0.419 = 0.632.

% latex table generated in R 4.5.1 by xtable 1.8-4 package
% Thu Sep 10 13:01:19 2026
\begin{table}[ht]
\centering
\begin{tabular}{rrrrr}
  \hline
 & Estimate & Std. Error & t value & Pr(>|t|) \\ 
  \hline
(Intercept) & 0.213 & 0.018 & 12.167 & 0.000 \\ 
  high\_income & 0.419 & 0.032 & 13.078 & 0.000 \\ 
   \hline
\end{tabular}
\caption{Coefficients of naive OLS regression} 
\label{tab:naive-regression}
\end{table}


That means that the naive OLS regression overestimates the propensity of high-income voters to vote in favour of the reform and underestimates the propensity of low-income voters to vote in favour of the reform, showing a more polarized distribution of votes conditional on income. 

That can be explained by the fact that the naive OLS attributes higher percentages of votes in favour in Northern precincts only to their higher share of high-income voters, while in reality the higher percentages of votes in favour are also driven by the fact that low-income voters in Northern precincts have a higher propensity to vote in favour of the reform than low-income voters in other regions.

\subsubsection{\texttt{seine} estimates}

When we move to semiparametric estimation performed by \texttt{seine}, we obtain the results presented in Tab. \ref{tab:ei-estimates}, which are much closer to the true values.

% latex table generated in R 4.5.1 by xtable 1.8-4 package
% Tue Aug 18 12:42:22 2026
\begin{table}[ht]
\centering
\begin{tabular}{llrrrr}
  \hline
predictor & outcome & estimate & std.error & conf.low & conf.high \\ 
  \hline
high\_income & outcome & 0.50 & 0.02 & 0.47 & 0.54 \\ 
  low\_income & outcome & 0.34 & 0.01 & 0.32 & 0.36 \\ 
   \hline
\end{tabular}
\caption{Semiparametric ecological inference estimates} 
\label{tab:ei-estimates}
\end{table}



%I may note that in this case, f(Z) was linear: 

%$f(Z_g) = 
%\begin{pmatrix}
%    0 & -0.3 \\
%    0.5 & 0.1
%\end{pmatrix} Z_g' + \begin{pmatrix}
%    0.6 \\
%    0.2
%\end{pmatrix}$

%Therefore, if we had correctly specified the model, we would have obtained the same results with a parametric estimation.

%In fact, $\beta_H = 0.6 + \begin{pmatrix}
%   0 & -0.3
%\end{pmatrix} Z_g'$ and $\beta_L = 0.2 + \begin{pmatrix}
%   0.5 & -0.1
%\end{pmatrix} Z_g'$.

%With 0.6, 0.2, 0, -0.3, 0.5, -0.1 quantities to be estimated that can be denoted by $\gamma_0^H$, $\gamma_0^L$, $\gamma_1^H$, $\gamma_1^L$, $\gamma_2^H$, $\gamma_2^L$.

%$\overline{Y_g}  = (\gamma_0^H+\gamma_1^H*Z^{North}+\gamma_2^H*Z^{Centre})*\overline{X}_g^H + (\gamma_0^L+\gamma_1^L*Z^{North}+\gamma_2^L*Z^{Centre})*\overline{X}_g^L + \epsilon_g$

%We do not know if \texttt{seine} has directly understood and exploited this linear structure, but we know that its estimates were as good as it had correctly understood the relationship.

%In the Appendix I can show another example:
%I might suppose that Z_g is a continous variable, for example the share of immigrants within a precinct.
%Therefore, there exist no regional distinction, precincts differ by their shares of immigrants.
%However, we have a relationship, even if loose, between the share of immigrants and \overline{X}_g.
%We should procede as follows:
%extract \overline{X_g} for each geography from a random distribution (e.g. \overline{X_g} & \sim \text{TruncNormal}(a=0, b=1, \mu=0.5, \sigma=0.2))
%that means that precincts differ for their percentage of high-income and low-income voters.
%Z_g & \sim TruncNormal(a=0, b=1, \mu=1-\overline{X}_g, \sigma=0.2)
%that means that there is a loose association between share of high-income voters and share of immigrants (interpretation: immigrants are more likely to live in low-income areas)

%Then, we might assume that 
%\begin{align*}
%\beta_{gr}^{H} &\sim \text{TruncNormal}(a=0, b=1, \mu=\eta_{gr}^{H}, \sigma=0.05), \\
%\beta_{gr}^{L} &\sim \text{TruncNormal}(a=0, b=1, \mu=\eta_{gr}^{L}, \sigma=0.05).
%\end{align*}

%\begin{align*}
%\eta_{gr}^{H} = e^{Z_g-1}, \quad & \eta_{gr}^{L} = 2^{Z_g-1}, \\
%\end{align*}

%interpretation: both high-income and low-income voters tend to favor the reform more as the share of immigrants in their precinct increases, but the effect is stronger for low-income voters and the relationship is non-linear

%we should observe that a linear estimator is not apt, while \texttt{seine} is able to capture the non-linear relationship and provide good estimates of $\beta_H$ and $\beta_L$